\documentclass[11pt]{amsart}

\usepackage[T1]{fontenc}
\usepackage{lmodern}
\usepackage{microtype}
\usepackage{amsmath,amssymb,amsthm}
\usepackage[hidelinks]{hyperref}

\hypersetup{
  pdftitle={The Hessian Arrangement Requires Four Added Lines for Supersolvability}
}

\newtheorem{theorem}{Theorem}
\newtheorem{lemma}[theorem]{Lemma}

\newcommand{\PP}{\mathbb{P}}
\newcommand{\Sing}{\operatorname{Sing}}
\newcommand{\extSS}{\operatorname{extSS}}

\title{The Hessian Arrangement Requires Four Added Lines for Supersolvability}
\subjclass[2020]{52C35, 32S22}
\keywords{Hessian arrangement, line arrangement, supersolvable arrangement,
modular point, extension to supersolvability}
\date{}

\begin{document}

\begin{abstract}
Let $\mathcal H$ be the twelve-line Hessian arrangement in
$\PP^2_{\mathbb C}$. We prove that $\extSS(\mathcal H)=4$. Thus no
extension of $\mathcal H$ by one line is supersolvable. Since
$\mathcal H$ is free with exponents $(4,7)$ and has maximal
multiplicity $4$, this gives a negative answer to the
supersolvable-extension question of Dimca, K\"uhne, and Pokora.
Dimca had already exhibited the relevant exceptional collinearity;
the contribution here is the obstruction for every possible modular
point and the exact extension number.
\end{abstract}

\maketitle

\section{Statement}

Let $\mathcal A$ be a line arrangement in $\PP^2_{\mathbb C}$. A point
$p\in\Sing(\mathcal A)$ is \emph{modular} if the line $pq$ belongs to
$\mathcal A$ for every $q\in\Sing(\mathcal A)\setminus\{p\}$. The
arrangement is \emph{supersolvable} if it has a modular point. For line
arrangements, that is, in rank three, this is equivalent to
supersolvability of the intersection lattice $L(\mathcal A)$, which is
the definition used in~\cite[Definition~1.1]{Kabat}; see~\cite{OT}, and
compare~\cite{DKP}, where a line arrangement having a modular point is
called supersolvable. Following
Kabat~\cite[Definition~1.3]{Kabat}, define
\[
\extSS(\mathcal A)
 =\min\bigl\{|\mathcal B\setminus\mathcal A|:
 \mathcal A\subseteq\mathcal B,\ \mathcal B\text{ supersolvable}\bigr\}.
\]
For line arrangements this is equivalent to Kabat's definition by
supersolvable resolutions, since the added lines may be ordered
arbitrarily.

Let $\omega$ be a primitive cube root of unity. The Hessian arrangement
is
\begin{equation}\label{eq:Hessian}
\mathcal H:\quad
xyz\bigl((x^3+y^3+z^3)^3-27x^3y^3z^3\bigr)=0.
\end{equation}
Its twelve lines are
\[
x=0,\qquad y=0,\qquad z=0,
\qquad x+\omega^a y+\omega^b z=0
\quad(a,b\in\mathbb Z/3\mathbb Z).
\]
Dimca~\cite[Example~6.5(i)]{Dimca} records that $\mathcal H$ is free
with first exponent $4=m(\mathcal H)$; hence its exponents are $(4,7)$.

\begin{theorem}\label{thm:main}
The Hessian arrangement satisfies
\[
\extSS(\mathcal H)=4.
\]
\end{theorem}

We recall the question this answers. Let $\mathcal A$ be an arrangement
of $d$ lines with first exponent $d_1$, and let $p\in\mathcal A$ be a
point of multiplicity $m_p$ with
\[
m_p\le d_1+1<d-m_p+1 .
\]
Then~\cite[Theorem~1.12(2)]{Dimca} states that if $m_p=d_1$ and
$m_p=m(\mathcal A)$, the maximal multiplicity of points of
$\mathcal A$, then all the other multiple points of $\mathcal A$ that
are not connected to $p$ by lines of $\mathcal A$ lie on a single line
$L$. Assuming $d_1=m(\mathcal A)$, Dimca, K\"uhne, and Pokora ask
in~\cite[Remark~2.3]{DKP} ``whether the line $L$
in~\cite[Theorem~1.12(2)]{Dimca} can be always chosen such that
$\mathcal A\cup L$ is supersolvable''. For $\mathcal H$ one has $d=12$
and $d_1=4=m(\mathcal H)$, and for a
point $p$ of multiplicity $m_p=4$ the hypothesis above reads
$4\le5<9$; thus $\mathcal H$ is in the scope of the question.
Theorem~\ref{thm:main} gives a negative answer.

\section{Incidences}

The singular points of $\mathcal H$ are
\begin{align*}
\mathcal P
 &=\bigl\{(0:1:-\omega^j),(1:0:-\omega^j),(1:-\omega^j:0):
 j\in\mathbb Z/3\mathbb Z\bigr\},\\
\mathcal V
 &=\bigl\{(1:0:0),(0:1:0),(0:0:1)\bigr\}
   \cup
   \bigl\{(1:\omega^r:\omega^s):
   r,s\in\mathbb Z/3\mathbb Z\bigr\}.
\end{align*}
Direct substitution shows that the nine points of $\mathcal P$ have
multiplicity four, the twelve points of $\mathcal V$ have multiplicity
two, and each line of $\mathcal H$ contains three points of
$\mathcal P$ and two points of $\mathcal V$. Moreover,
\[
9\binom42+12\binom22=66=\binom{12}{2},
\]
so these are all the singular points.

\begin{lemma}\label{lem:incidence}
The following hold.
\begin{enumerate}
\item Every two points of $\mathcal P$ lie on a line of $\mathcal H$.
\item For each $p\in\mathcal P$, exactly four points of $\mathcal V$
are not joined to $p$ by a line of $\mathcal H$. Those four points are
collinear on a line not containing $p$.
\item Let $v\in\mathcal V$. The two lines of $\mathcal H$ through $v$
contain six points of $\mathcal P$. If $b_1,b_2,b_3$ are the remaining
points of $\mathcal P$, then each line $vb_i$ contains no point of
$\mathcal V$ other than $v$.
\end{enumerate}
\end{lemma}

\begin{proof}
Each line of $\mathcal H$ contains three points of $\mathcal P$, and
\[
12\binom32=\binom92.
\]
No pair is counted twice, since two distinct points determine a unique
projective line. This proves~(1).

Coordinate permutations and diagonal maps
\[
(x:y:z)\longmapsto(\omega^a x:\omega^b y:\omega^c z)
\]
preserve $\mathcal H$ and act transitively on $\mathcal P$. For~(2), it
therefore suffices to take $p=(1:-1:0)$. The four exceptional points are
\[
(0:0:1),\qquad (1:1:1),\qquad
(1:1:\omega),\qquad (1:1:\omega^2).
\]
They lie on $x-y=0$, which does not contain $p$.

The same symmetries have two orbits on $\mathcal V$: the coordinate
vertices and the points with all coordinates nonzero. For
$v=(0:0:1)$, the two old lines through $v$ are $x=0$ and $y=0$, and the
remaining points of $\mathcal P$ are
\[
b_j=(1:-\omega^j:0),\qquad j\in\mathbb Z/3\mathbb Z.
\]
Their joins with $v$ are $\omega^j x+y=0$. None contains another point
of $\mathcal V$: substituting $(1:0:0)$ gives $\omega^j\ne0$ and
substituting $(0:1:0)$ gives $1\ne0$, while at $(1:\omega^r:\omega^s)$
the equation would force $\omega^j+\omega^r=0$, which is impossible for
cube roots of unity.

For $v=(1:1:1)$, the remaining points of $\mathcal P$ are
\[
(0:1:-1),\qquad (1:0:-1),\qquad (1:-1:0),
\]
and their joins with $v$ are
\[
2x-y-z=0,\qquad x-2y+z=0,\qquad x+y-2z=0.
\]
The first contains no coordinate vertex. At a point
$(1:\omega^r:\omega^s)$, its equation gives
$\omega^r+\omega^s=2$, hence $r=s=0$. The other two cases follow by
permuting the coordinates. This proves~(3).
\end{proof}

\section{Proof}

Let $\mathcal B\supseteq\mathcal H$ be supersolvable, set
$k=|\mathcal B\setminus\mathcal H|$, and let $X$ be a modular point of
$\mathcal B$.

Suppose first that $X\in\mathcal P$. By
Lemma~\ref{lem:incidence}(2), there are four points
$v_1,\dots,v_4\in\mathcal V$ not joined to $X$ by old lines. Modularity
forces the four joins $Xv_i$ to belong to $\mathcal B$. They are
outside $\mathcal H$ and are distinct: the points $v_i$ lie on a line
not containing $X$, so no line through $X$ contains two of them. Thus
$k\ge4$.

Suppose next that $X\in\mathcal V$. The two old lines through $X$
contain six points of $\mathcal P$. The remaining three points require
three distinct added lines through $X$, because a line outside
$\mathcal H$ contains at most one point of $\mathcal P$ by
Lemma~\ref{lem:incidence}(1). By Lemma~\ref{lem:incidence}(3), none of
these three lines contains a point of $\mathcal V\setminus\{X\}$. The
union of the two old lines through $X$ contains only three points of
$\mathcal V$, including $X$. Hence some point of $\mathcal V$ requires
a fourth added line through $X$, and again $k\ge4$.

Finally, suppose that $X\notin\Sing(\mathcal H)$. Every point of
$\mathcal P$ remains singular in $\mathcal B$, so modularity requires
its join with $X$ to belong to $\mathcal B$. At most one old line
passes through $X$, and it contains at most three points of
$\mathcal P$. Every added line through $X$ contains at most one point
of $\mathcal P$. If $k\le3$, at most six of the nine points of
$\mathcal P$ can be joined to $X$, a contradiction. Therefore
\[
\extSS(\mathcal H)\ge4.
\]
The three cases give lower bounds only; the individual counts are not
claimed to be sharp, and $k\ge4$ is all that is needed.

For the reverse inequality, take $p=(1:-1:0)$ and set
\[
\mathcal B=\mathcal H
\cup\{x+y=0\}
\cup\{x+y-2\omega^jz=0:j\in\mathbb Z/3\mathbb Z\}.
\]
These are precisely the four joins from $p$ to the exceptional points
in Lemma~\ref{lem:incidence}(2). Every old singular point is joined to
$p$ by a line of $\mathcal B$. If
$q\in\Sing(\mathcal B)\setminus\Sing(\mathcal H)$, then the
intersection defining $q$ involves an added line; hence $q$ lies on an
added line through $p$. Thus $p$ is modular in $\mathcal B$, and
\[
\extSS(\mathcal H)\le4.
\]
This proves Theorem~\ref{thm:main}.

For $p=(1:-1:0)$, the four singular points not joined to $p$ by lines
of $\mathcal H$ are distinct and determine the unique line $x-y=0$.
Thus this is the line supplied by~\cite[Theorem~1.12(2)]{Dimca}.
Since no one-line extension of $\mathcal H$ is supersolvable, this
proves the claimed negative answer to~\cite[Remark~2.3]{DKP}.

\begin{thebibliography}{9}

\bibitem{Dimca}
A.~Dimca,
\emph{On the module of derivations of a line arrangement},
to appear in Publ. Mat.;
\href{https://arxiv.org/abs/2503.01624v7}{arXiv:2503.01624v7}.

\bibitem{DKP}
A.~Dimca, L.~K\"uhne, and P.~Pokora,
\emph{Free line arrangements with low maximal multiplicity},
J. Algebraic Combin. \textbf{63} (2026), Art.~45,
\href{https://doi.org/10.1007/s10801-026-01533-8}
{doi:10.1007/s10801-026-01533-8};
\href{https://arxiv.org/abs/2505.01733v4}{arXiv:2505.01733v4}.
Numbering and the quotation in the text follow the preprint.

\bibitem{Kabat}
J.~Kabat,
\emph{Supersolvable resolutions of line arrangements},
Hokkaido Math. J. \textbf{52} (2023), no.~2, 301--313,
\href{https://doi.org/10.14492/hokmj/2021-540}
{doi:10.14492/hokmj/2021-540};
\href{https://arxiv.org/abs/2201.04856v1}{arXiv:2201.04856v1}.
Numbering of definitions follows the preprint.

\bibitem{OT}
P.~Orlik and H.~Terao,
\emph{Arrangements of hyperplanes},
Springer-Verlag, Berlin, 1992.

\end{thebibliography}

\end{document}